\section{Families of directions}
\label{dir:section}

The construction uses many linear images of the same parameter space.  The
directions defining each image must approximate a Gaussian transform, while
most images must have small projections in any prescribed direction.
Sampling Fourier features and controlling them on finite nets is a standard
route to uniform kernel approximation~\cite{RR07}. Here the same samples
must also control derivative profiles and satisfy a majority-chart
guarantee in every direction.

\begin{lemma}\label{lem:directions}
There is an absolute constant $C$ such that, for every sufficiently large
integer $m$, the following holds.  Set $N=m^6$ and $S=m^3$.  There are vectors
$g_{s,i}\in\R^m$, indexed by $1\le s\le S$ and $1\le i\le N$, with these
properties:
\begin{enumerate}
\item For every $s,i$,
\begin{equation}\label{dir:norms}
 \frac{\sqrt m}{2}\le |g_{s,i}|\le 2\sqrt m.
\end{equation}
Within each family $s$, the vectors have pairwise nonparallel directions.
Every eigenvalue of $N^{-1}\sum_{i=1}^N g_{s,i}g_{s,i}^{T}$ belongs to
$[1/2,2]$.
\item For every $s$,
\begin{equation}\label{dir:gaussian-approximation}
 \sup_{|w|\le m^{42}}
 \left|\frac1N\sum_{i=1}^N
 g_{s,i}\sin\langle g_{s,i},w\rangle
       -w e^{-|w|^2/2}\right|
 \le \frac{2}{\sqrt m}.
\end{equation}
\item For every unit vector $v\in\R^m$, at least $S/2$ of the indices $s$
satisfy
\begin{equation}\label{dir:good-charts}
 \max_{1\le i\le N}|\langle v,g_{s,i}\rangle|
 \le C\sqrt{\log m}.
\end{equation}
\end{enumerate}
\end{lemma}

Independent standard Gaussian vectors give such a family. The frequency
net needed in the proof has $\exp(O(m\log m))$ points, whereas the
failure probability for each scalar empirical estimate is at most
$\exp(-cm^4)$. A separate sphere net and independence across charts give
the last assertion simultaneously for every unit direction. The full
probabilistic argument appears in Appendix~\ref{app:directions}.


Fix one such family for the remainder of the proof. In particular, its
every-direction guarantee is available before any contraction is chosen.
Define
\begin{equation}\label{dir:chart-maps}
 u_{s,i}=\frac{g_{s,i}}{\sqrt m},
 \qquad
 J_s:\R^m\longrightarrow\R^N,
 \qquad
 (J_s\theta)_i=\langle u_{s,i},\theta\rangle.
\end{equation}
The eigenvalues of $J_s^TJ_s$ lie in $[N/(2m),2N/m]$, so $J_s$ has full
column rank.  Its left pseudoinverse
$J_s^+=(J_s^TJ_s)^{-1}J_s^T$ satisfies $J_s^+J_s=I_m$, and $J_sJ_s^+$ is the
orthogonal projection onto the range of $J_s$.  We will use
\begin{equation}\label{eq:pseudoinverse}
 \|J_s\|_{\op}\le\sqrt{\frac{2N}{m}},
 \qquad
 \|J_s^+\|_{\op}\le2\sqrt{\frac mN},
 \qquad
 |J_s^+e_i|\le\frac{4m}{N}\quad(1\le i\le N),
\end{equation}
where $e_i$ is the $i$th standard coordinate vector of $\R^N$.  The final
bound follows from $J_s^+e_i=(J_s^TJ_s)^{-1}u_{s,i}$,
$\|(J_s^TJ_s)^{-1}\|_{\op}\le2m/N$, and $|u_{s,i}|\le2$.
